\section{树、图、BFS 和 DFS 章正文案例推导补强}

这一节补齐本章正文出现过的 LeetCode 案例推导。阅读顺序统一按“题目说明、样例入口、暴力基线、暴力过程、重复工作、优化条件、相邻状态、状态复用、指针和字段、代码落点、正确性、边界实验、复杂度变化、迁移追问”展开；目的不是新增题量，而是把正文中的案例都讲成可复述、可证明、可迁移的解题链。

\subsection{LeetCode 102 二叉树层序遍历}

\textbf{题目说明。} LeetCode 102 二叉树层序遍历 的题面入口要先还原成具体任务：队列在每轮开始时保存当前层所有节点。

\textbf{样例入口。} 拿根 3、左 9、右 20 手算，队列第一轮只有 3，第二轮才是 9 和 20。若不固定本层 size，新入队的下一层节点会混进当前层。

\textbf{暴力基线。} 慢版本可以先给每个节点计算深度，再按深度分组；或者每层都从根重新扫描一遍收集该层节点。

\textbf{暴力过程。} 以根 3、左右孩子 9 和 20 为例，第一层只有 3，第二层是 9、20。若队列不固定当前层长度，20 的孩子可能混进同一层。

\textbf{重复工作。} 题面模型：队列在每轮开始时保存当前层所有节点。暴力版的慢点要落到本题：浪费在重复从根找层，或者层边界不清导致混层。队列天然按访问顺序保存下一批节点，只要固定本轮队列大小就是当前层。后面的优化只消掉这类重复工作，不能改变答案集合。

\textbf{优化条件。} 本题的关键观察是：记录 level size 固定层边界，避免下一层节点混入当前层。这句话必须能翻译成可维护状态；如果题目条件改变后观察不再成立，就不能继续套原来的写法。

\textbf{相邻状态。} 规律是每轮开始记录队列长度，这个长度就是当前层边界。把这个特例继续拆开看：节点向父亲返回的量和全局答案通常不是同一个量。先写叶子或空节点的返回值，再写父节点如何合并左右孩子，递归式才有边界基础。

\textbf{状态复用。} BFS 队列保存待处理节点。每轮先读取 level\_size，只弹出这 level\_size 个节点并收集值，同时把它们的孩子加入队尾。手算时只改被新输入影响的那一小部分，而不是回头重算完整候选。

\textbf{指针和字段。} 状态字段要能说清：queue 是待访问节点队列，level\_size 是当前层节点数，level\_values 是本层输出。手算时按这个协议跟踪：队列初始 [3]，level\_size=1，弹出 3 后压入 9 和 20；下一轮 level\_size=2，只处理 9 和 20。

\textbf{代码落点。} 关键是 for 循环次数使用进入本层前的 queue.size()，不要在循环条件里直接用动态变化的 queue.size()。关键代码行必须能逐句翻译成“旧状态怎样变成新状态”，否则就只是模板。

\textbf{正确性。} BFS 入队顺序保证父层节点先于子层节点；固定 level\_size 后，本轮处理的恰好是同一深度的全部节点。

\textbf{边界实验。} 先用空树、只有根、最后一层不满、左右顺序做反例检查。实验时覆盖空树、只有根、最后一层不满、左右顺序，尤其关注空节点、单节点、链式树和全负值。递归版本和显式栈版本可以互相对拍。

\textbf{复杂度变化。} 重复扫描每层最坏 O(nh)，BFS 每节点入队出队一次 O(n)，队列空间 O(width)。

\textbf{迁移追问。} 如果题目改成“锯齿形层序只是在每层输出顺序上增加方向控制”，先判断原状态是否仍然覆盖新答案集合，再决定是微调边界、改 value 语义，还是换算法。

\subsection{LeetCode 200 岛屿数量}

\textbf{题目说明。} LeetCode 200 岛屿数量 的题面入口要先还原成具体任务：陆地格子是节点，四方向相邻是边，岛屿是连通块。

\textbf{样例入口。} 拿一个 2x2 小岛手算，从一个陆地出发能沿上下左右标记同一块岛。若不标记访问，会重复计数；若对角线也连通，会误合并。

\textbf{暴力基线。} 暴力可以对每个陆地格子单独搜索它所属的连通块，再用集合去重；这样同一块岛会被重复搜索很多次。

\textbf{暴力过程。} 以 2x2 全是陆地为例，从左上搜索会访问四个格子；如果之后从右上又重新搜索，就把同一座岛重复数了一遍。

\textbf{重复工作。} 题面模型：陆地格子是节点，四方向相邻是边，岛屿是连通块。暴力版的慢点要落到本题：浪费在已确认属于某个岛的陆地没有被标记。每块岛只需要在第一次遇到时完整淹没或标记。后面的优化只消掉这类重复工作，不能改变答案集合。

\textbf{优化条件。} 本题的关键观察是：扫描遇到未访问陆地就启动搜索并标记整块。这句话必须能翻译成可维护状态；如果题目条件改变后观察不再成立，就不能继续套原来的写法。

\textbf{相邻状态。} 规律是每发现一个未访问陆地，答案加一，并用 DFS/BFS 标记与它四连通的全部陆地。把这个特例继续拆开看：状态节点不一定等于原始位置，还可能包含钥匙、剩余资源、时间或访问集合。visited 只能合并未来能力完全相同的状态，否则会把合法路径剪掉。

\textbf{状态复用。} 扫描网格。遇到未访问陆地，答案加一，然后 BFS/DFS 把与它四方向连通的全部陆地标记为已访问。手算时只改被新输入影响的那一小部分，而不是回头重算完整候选。

\textbf{指针和字段。} 状态字段要能说清：row/col 是当前格子，visited 或原地改写表示已经归属某座岛，directions 是四方向邻居。手算时按这个协议跟踪：从左上陆地入队，依次扩展右边和下边陆地，最终整块连通区域都标记；扫描到这些格子时直接跳过。

\textbf{代码落点。} 关键是只看上下左右，不看对角线。边界判断集中写，访问标记在入队时设置，避免重复入队。关键代码行必须能逐句翻译成“旧状态怎样变成新状态”，否则就只是模板。

\textbf{正确性。} 一次搜索会访问且只访问与起点四连通的全部陆地；外层每次启动搜索都对应一块此前未计数的新岛。

\textbf{边界实验。} 先用空网格、全水、全陆、斜对角不连通、是否允许修改输入做反例检查。实验时覆盖空网格、全水、全陆、斜对角不连通、是否允许修改输入，并打印入队时的完整状态。若同一位置但资源不同的状态被合并，很多图题会出现隐蔽错误。

\textbf{复杂度变化。} 重复搜索最坏 O(\ensuremath{(mn)^{2}})，标记搜索 O(mn) 时间，空间 O(mn) 或递归/队列大小。

\textbf{迁移追问。} 如果题目改成“也可用并查集合并相邻陆地，适合动态岛屿扩展”，先判断原状态是否仍然覆盖新答案集合，再决定是微调边界、改 value 语义，还是换算法。

\subsection{LeetCode 207 课程表}

\textbf{题目说明。} LeetCode 207 课程表 的题面入口要先还原成具体任务：课程是节点，先修关系是有向边，问题是判断是否有环。

\textbf{样例入口。} 拿课程 0->1 手算，入度为 0 的课程 0 先入队，学完后 1 的入度降为 0；若有 0->1 和 1->0，队列一开始为空。

\textbf{暴力基线。} 暴力可以不断扫描所有课程，找当前先修都已满足的课程；若一轮没有新课程可学，就判断有环。

\textbf{暴力过程。} 以 0->1 为例，课程 0 入度为 0，先学 0，随后 1 的入度变 0。若有 0->1 和 1->0，一开始没有入度为 0 的课。

\textbf{重复工作。} 题面模型：课程是节点，先修关系是有向边，问题是判断是否有环。暴力版的慢点要落到本题：浪费在每轮都重新检查所有课程的先修。入度正好表示还剩多少未满足依赖，学完一门课只影响它指向的后继。后面的优化只消掉这类重复工作，不能改变答案集合。

\textbf{优化条件。} 本题的关键观察是：入度拓扑排序不断学习当前无依赖课程。这句话必须能翻译成可维护状态；如果题目条件改变后观察不再成立，就不能继续套原来的写法。

\textbf{相邻状态。} 规律是拓扑排序用入度为 0 的节点推进，处理数量小于课程数说明有环。把这个特例继续拆开看：状态节点不一定等于原始位置，还可能包含钥匙、剩余资源、时间或访问集合。visited 只能合并未来能力完全相同的状态，否则会把合法路径剪掉。

\textbf{状态复用。} 构建邻接表和 indegree。把所有入度为 0 的课程入队，弹出课程时让后继入度减一，减到 0 就入队。手算时只改被新输入影响的那一小部分，而不是回头重算完整候选。

\textbf{指针和字段。} 状态字段要能说清：indegree[v] 是课程 v 尚未满足的先修数量，queue 保存当前可学习课程，visited\_count 是已处理数量。手算时按这个协议跟踪：0 出队后，遍历它的后继 1，indegree[1] 从 1 变 0，1 入队；最后处理数量等于课程数，说明无环。

\textbf{代码落点。} 关键是边方向要统一：prerequisite [a, b] 表示 b -> a。处理数量小于 numCourses 时返回 false。关键代码行必须能逐句翻译成“旧状态怎样变成新状态”，否则就只是模板。

\textbf{正确性。} 入度为 0 的节点没有未满足前置条件，可以安全学习；有向环中的节点永远不会全部降到入度 0。

\textbf{边界实验。} 先用边方向、孤立课程、环、自依赖、重复边做反例检查。实验时覆盖边方向、孤立课程、环、自依赖、重复边，并打印入队时的完整状态。若同一位置但资源不同的状态被合并，很多图题会出现隐蔽错误。

\textbf{复杂度变化。} 反复扫描最坏 O(VE)，拓扑排序 O(V+E) 时间和空间。

\textbf{迁移追问。} 如果题目改成“DFS 三色标记也能判环，但要明确访问状态”，先判断原状态是否仍然覆盖新答案集合，再决定是微调边界、改 value 语义，还是换算法。

\subsection{LeetCode 127 单词接龙}

\textbf{题目说明。} LeetCode 127 单词接龙 的题面入口要先还原成具体任务：每个单词是节点，一次修改一个字符形成边。

\textbf{样例入口。} 拿 hit -> hot -> dot -> dog -> cog 手算，每次只能改一个字符，所以 BFS 第一层是 hot，第二层才到 dot 或 lot。

\textbf{暴力基线。} 暴力可以把每个单词和所有其他单词比较一遍，差一个字符就连边，再从 beginWord 做 BFS。

\textbf{暴力过程。} 以 hit 到 cog 为例，hit 只能一步到 hot；hot 下一层到 dot 或 lot。第一次到 cog 的层数就是最短转换长度。

\textbf{重复工作。} 题面模型：每个单词是节点，一次修改一个字符形成边。暴力版的慢点要落到本题：浪费在找邻居时每次扫描整个词表。单词长度固定时，可以把 h*t、*ot 这类通配模式作为桶，快速找到只差一位的邻居。后面的优化只消掉这类重复工作，不能改变答案集合。

\textbf{优化条件。} 本题的关键观察是：通配模式桶加速邻居查询，BFS 第一次到达终点就是最短转换长度。这句话必须能翻译成可维护状态；如果题目条件改变后观察不再成立，就不能继续套原来的写法。

\textbf{相邻状态。} 规律是单词作为节点，通配桶如 h*t 把邻居快速找出；第一次到达终点就是最短转换长度。把这个特例继续拆开看：状态节点不一定等于原始位置，还可能包含钥匙、剩余资源、时间或访问集合。visited 只能合并未来能力完全相同的状态，否则会把合法路径剪掉。

\textbf{状态复用。} 预处理 pattern -> words 的映射，BFS 队列保存当前单词和步数，visited 防止重复入队。手算时只改被新输入影响的那一小部分，而不是回头重算完整候选。

\textbf{指针和字段。} 状态字段要能说清：word 是当前单词，pattern 是替换某一位后的通配键，distance 是转换层数。手算时按这个协议跟踪：hit 生成 *it、h*t、hi*，其中 h*t 桶找到 hot；hot 再通过 *ot 找到 dot 和 lot。

\textbf{代码落点。} 关键是访问标记在入队时设置。桶用完后可以清空，避免不同节点重复展开同一批邻居。关键代码行必须能逐句翻译成“旧状态怎样变成新状态”，否则就只是模板。

\textbf{正确性。} 每条转换边权都为 1，BFS 按层扩展；第一次到达 endWord 时经过的边数最少。

\textbf{边界实验。} 先用终点不在词表、起点和终点相同、桶重复扫描、单词长度一致做反例检查。实验时覆盖终点不在词表、起点和终点相同、桶重复扫描、单词长度一致，并打印入队时的完整状态。若同一位置但资源不同的状态被合并，很多图题会出现隐蔽错误。

\textbf{复杂度变化。} 两两建图 O(\ensuremath{N^{2}}L)，通配桶建图和 BFS 约 O(\ensuremath{NL^{2}}) 或按桶规模计，空间 O(NL)。

\textbf{迁移追问。} 如果题目改成“双向 BFS 可以进一步降低搜索空间”，先判断原状态是否仍然覆盖新答案集合，再决定是微调边界、改 value 语义，还是换算法。

\subsection{LeetCode 864 获取所有钥匙的最短路径}

\textbf{题目说明。} LeetCode 864 获取所有钥匙的最短路径 的题面入口要先还原成具体任务：位置相同但钥匙集合不同，未来能开的门不同。

\textbf{样例入口。} 拿网格 @.a / \#.\# / b.A 手算，走到 a 后状态不是位置变了而是钥匙集合多了一把；同一位置带不同钥匙未来能力不同。

\textbf{暴力基线。} 慢版本先按题意画出完整状态图，用普通搜索跑通可达性。这个过程用于确认位置相同但钥匙集合不同，未来能开的门不同的节点和边，之后再讨论访问标记、层次或代价顺序。放到本题就是：先按“位置相同但钥匙集合不同，未来能开的门不同”完整试慢版本；样例里已经能看到“规律是 BFS 的 visited 必须按 (row, col, keymask) 记录，拿齐所有钥匙第一次出队就是最短步数”，所以优化前必须先能说清慢版本枚举了哪些候选。

\textbf{暴力过程。} 拿网格 @.a / \#.\# / b.A 手算，走到 a 后状态不是位置变了而是钥匙集合多了一把；同一位置带不同钥匙未来能力不同。

\textbf{重复工作。} 题面模型：位置相同但钥匙集合不同，未来能开的门不同。暴力版的慢点要落到本题：慢版本会围绕“位置相同但钥匙集合不同，未来能开的门不同”反复展开同一批状态。样例规律是“规律是 BFS 的 visited 必须按 (row, col, keymask) 记录，拿齐所有钥匙第一次出队就是最短步数”，说明必须把节点、边、访问标记和资源维度一次建清；同一状态不应重复入队，不同资源状态也不能误合并。后面的优化只消掉这类重复工作，不能改变答案集合。

\textbf{优化条件。} 本题的关键观察是：BFS 状态包含行、列和钥匙掩码，访问标记也必须包含掩码。这句话必须能翻译成可维护状态；如果题目条件改变后观察不再成立，就不能继续套原来的写法。

\textbf{相邻状态。} 规律是 BFS 的 visited 必须按 (row, col, keymask) 记录，拿齐所有钥匙第一次出队就是最短步数。把这个特例继续拆开看：状态节点不一定等于原始位置，还可能包含钥匙、剩余资源、时间或访问集合。visited 只能合并未来能力完全相同的状态，否则会把合法路径剪掉。

\textbf{状态复用。} 把重复搜索压成访问标记、距离数组或拓扑状态；邻居生成也要从全量扫描变成结构化枚举。本题落点是：规律是 BFS 的 visited 必须按 (row, col, keymask) 记录，拿齐所有钥匙第一次出队就是最短步数。手算时只改被新输入影响的那一小部分，而不是回头重算完整候选。

\textbf{指针和字段。} 状态字段要能说清：状态编号、邻接生成方式、访问标记、距离或层数、队列内容；若状态含资源，visited 不能只按位置记录。手算时按这个协议跟踪：拿网格 @.a / \#.\# / b.A 手算，走到 a 后状态不是位置变了而是钥匙集合多了一把；同一位置带不同钥匙未来能力不同。然后按协议复查：手算时画队列或栈，状态必须包含题目需要的所有资源。入队时标记还是出队时标记，要和去重语义一致。

\textbf{代码落点。} 入队时标记还是出队时标记必须一致；方向数组集中定义；多源 BFS 要一次性把所有源点放入初始队列。关键代码行必须能逐句翻译成“旧状态怎样变成新状态”，否则就只是模板。

\textbf{正确性。} 证明从可达性开始：搜索覆盖所有合法边能到达的状态，访问标记只删除重复状态而不删除不同未来能力。

\textbf{边界实验。} 先用起点旁边有门、钥匙顺序、不可达、钥匙数量决定掩码大小做反例检查。实验时覆盖起点旁边有门、钥匙顺序、不可达、钥匙数量决定掩码大小，并打印入队时的完整状态。若同一位置但资源不同的状态被合并，很多图题会出现隐蔽错误。

\textbf{复杂度变化。} 暴力重复从多个起点搜索会反复遍历图；正确建模和访问标记后，BFS 或 DFS 通常按节点数加边数计算，状态图题则按完整状态数计算。

\textbf{迁移追问。} 如果题目改成“这是状态图维度不能只按位置去重的典型题”，先判断原状态是否仍然覆盖新答案集合，再决定是微调边界、改 value 语义，还是换算法。

\subsection{LeetCode 104 二叉树最大深度}

\textbf{题目说明。} LeetCode 104 二叉树最大深度 的题面入口要先还原成具体任务：深度可以看成层数，也可以作为栈中携带的状态。

\textbf{样例入口。} 拿链式树 1->2->3 手算，最大深度是 3；空树是 0。递归返回左右子树深度的较大值加一，显式栈则把节点和当前深度一起压栈。

\textbf{暴力基线。} 慢版本在每个节点重新扫描子树或路径，先求出局部答案。重复扫描暴露了真正状态：每个节点只需要从孩子或父亲那里拿到少量信息。放到本题就是：先按“深度可以看成层数，也可以作为栈中携带的状态”完整试慢版本；样例里已经能看到“规律是状态必须携带深度，不能只记录访问过哪些节点”，所以优化前必须先能说清慢版本枚举了哪些候选。

\textbf{暴力过程。} 拿链式树 1->2->3 手算，最大深度是 3；空树是 0。递归返回左右子树深度的较大值加一，显式栈则把节点和当前深度一起压栈。

\textbf{重复工作。} 题面模型：深度可以看成层数，也可以作为栈中携带的状态。暴力版的慢点要落到本题：慢版本会在树上重复确认“深度可以看成层数，也可以作为栈中携带的状态”。样例规律是“规律是状态必须携带深度，不能只记录访问过哪些节点”，说明父子之间真正需要传递的是高度、上下界、命中状态、层号或两类选择值，而不是反复扫描整棵子树。后面的优化只消掉这类重复工作，不能改变答案集合。

\textbf{优化条件。} 本题的关键观察是：显式栈保存节点和深度，避免极端链式树递归栈过深。这句话必须能翻译成可维护状态；如果题目条件改变后观察不再成立，就不能继续套原来的写法。

\textbf{相邻状态。} 规律是状态必须携带深度，不能只记录访问过哪些节点。把这个特例继续拆开看：节点向父亲返回的量和全局答案通常不是同一个量。先写叶子或空节点的返回值，再写父节点如何合并左右孩子，递归式才有边界基础。

\textbf{状态复用。} 把对子树的重复扫描压成返回值或队列状态；每个节点只合并孩子给出的稳定信息。本题落点是：规律是状态必须携带深度，不能只记录访问过哪些节点。手算时只改被新输入影响的那一小部分，而不是回头重算完整候选。

\textbf{指针和字段。} 状态字段要能说清：节点指针、传入约束或返回值、当前答案、层号或路径状态；空节点的返回值必须和状态定义匹配。手算时按这个协议跟踪：拿链式树 1->2->3 手算，最大深度是 3；空树是 0。递归返回左右子树深度的较大值加一，显式栈则把节点和当前深度一起压栈。然后按协议复查：手算时从叶子开始写返回值，或者从根开始写传入约束。不要只写访问顺序，要写每条边上传递的信息。

\textbf{代码落点。} 递归版本先处理空节点；迭代版本的栈元素要带完整状态；全负值、重复值和极端深树要单独测。关键代码行必须能逐句翻译成“旧状态怎样变成新状态”，否则就只是模板。

\textbf{正确性。} 证明从子树归纳或层序不变量开始：孩子状态正确时父节点合并正确；若用队列，当前轮队列必须正好保存当前层节点。

\textbf{边界实验。} 先用空树返回零、根深度是一、链式树做反例检查。实验时覆盖空树返回零、根深度是一、链式树，尤其关注空节点、单节点、链式树和全负值。递归版本和显式栈版本可以互相对拍。

\textbf{复杂度变化。} 暴力在每个节点重复扫描子树或反复寻找层边界可能退化到平方级；后序汇总、显式栈或层序遍历让每个节点处理常数次，通常是线性时间。

\textbf{迁移追问。} 如果题目改成“最小深度不能简单取左右最小，因为空子树不构成叶子路径”，先判断原状态是否仍然覆盖新答案集合，再决定是微调边界、改 value 语义，还是换算法。

\subsection{LeetCode 98 验证二叉搜索树}

\textbf{题目说明。} LeetCode 98 验证二叉搜索树 的题面入口要先还原成具体任务：BST 约束是全局上下界，不是只比较父子节点。

\textbf{样例入口。} 拿 [5, 1, 4, null, null, 3, 6] 手算，节点 3 小于根 5 却在右子树里，只比较父子 4 和 3 会误判。

\textbf{暴力基线。} 错误暴力常只检查当前节点大于左孩子、小于右孩子；正确慢版本要对每个节点扫描整棵左子树最大值和右子树最小值。

\textbf{暴力过程。} 以 [5, 1, 4, null, null, 3, 6] 为例，节点 4 的左孩子 3 小于 4、右孩子 6 大于 4，看局部没问题；但 3 在根 5 的右子树里，小于 5，所以整棵树不是 BST。

\textbf{重复工作。} 题面模型：BST 约束是全局上下界，不是只比较父子节点。暴力版的慢点要落到本题：浪费在每个节点重复扫描子树求上下界。真正需要传递的是祖先给当前子树限定的取值范围。后面的优化只消掉这类重复工作，不能改变答案集合。

\textbf{优化条件。} 本题的关键观察是：中序严格递增或显式传递上下界都可以；中序版本要保存前一个访问值。这句话必须能翻译成可维护状态；如果题目条件改变后观察不再成立，就不能继续套原来的写法。

\textbf{相邻状态。} 规律是 BST 约束来自整条祖先上下界，或中序序列必须严格递增；重复值和 int 边界要按题意单独处理。把这个特例继续拆开看：节点向父亲返回的量和全局答案通常不是同一个量。先写叶子或空节点的返回值，再写父节点如何合并左右孩子，递归式才有边界基础。

\textbf{状态复用。} 自顶向下传入 lower 和 upper，当前节点值必须满足 lower<val<upper；去左子树时 upper 变成当前值，去右子树时 lower 变成当前值。手算时只改被新输入影响的那一小部分，而不是回头重算完整候选。

\textbf{指针和字段。} 状态字段要能说清：node 是当前节点，lower/upper 是祖先约束；空节点返回合法。手算时按这个协议跟踪：到样例中右子树节点 3 时，它的 lower 应该是 5、upper 是 4 或来自路径约束，3 不满足根给的下界。

\textbf{代码落点。} 关键是用 long long 或 optional 边界，避免节点值等于 int 最小/最大时哨兵冲突；重复值是否允许按题意处理，本题严格不允许。关键代码行必须能逐句翻译成“旧状态怎样变成新状态”，否则就只是模板。

\textbf{正确性。} BST 要求左子树所有值小于根、右子树所有值大于根；上下界把所有祖先约束都带到当前节点，覆盖了非局部违规。

\textbf{边界实验。} 先用重复值、int 最小最大值、局部合法但全局非法做反例检查。实验时覆盖重复值、int 最小最大值、局部合法但全局非法，尤其关注空节点、单节点、链式树和全负值。递归版本和显式栈版本可以互相对拍。

\textbf{复杂度变化。} 扫描子树慢版本最坏 O(\ensuremath{n^{2}})，上下界 DFS/BFS 每节点一次 O(n)，空间 O(h) 或显式栈 O(h)。

\textbf{迁移追问。} 如果题目改成“若允许重复值，需要明确重复放左还是放右”，先判断原状态是否仍然覆盖新答案集合，再决定是微调边界、改 value 语义，还是换算法。

\subsection{LeetCode 236 二叉树最近公共祖先}

\textbf{题目说明。} LeetCode 236 二叉树最近公共祖先 的题面入口要先还原成具体任务：普通树没有有序性，只能依靠祖先关系。

\textbf{样例入口。} 拿普通二叉树根 3，p=5、q=1 手算，左右子树分别找到一个目标，根 3 就是最近公共祖先；若两个目标都在左子树，答案应由左子树返回。

\textbf{暴力基线。} 慢版本在每个节点重新扫描子树或路径，先求出局部答案。重复扫描暴露了真正状态：每个节点只需要从孩子或父亲那里拿到少量信息。放到本题就是：先按“普通树没有有序性，只能依靠祖先关系”完整试慢版本；样例里已经能看到“规律是后序返回是否找到目标或祖先，不能用 BST 的大小关系”，所以优化前必须先能说清慢版本枚举了哪些候选。

\textbf{暴力过程。} 拿普通二叉树根 3，p=5、q=1 手算，左右子树分别找到一个目标，根 3 就是最近公共祖先；若两个目标都在左子树，答案应由左子树返回。

\textbf{重复工作。} 题面模型：普通树没有有序性，只能依靠祖先关系。暴力版的慢点要落到本题：慢版本会在树上重复确认“普通树没有有序性，只能依靠祖先关系”。样例规律是“规律是后序返回是否找到目标或祖先，不能用 BST 的大小关系”，说明父子之间真正需要传递的是高度、上下界、命中状态、层号或两类选择值，而不是反复扫描整棵子树。后面的优化只消掉这类重复工作，不能改变答案集合。

\textbf{优化条件。} 本题的关键观察是：父指针集合或后序返回目标命中状态都可行。这句话必须能翻译成可维护状态；如果题目条件改变后观察不再成立，就不能继续套原来的写法。

\textbf{相邻状态。} 规律是后序返回是否找到目标或祖先，不能用 BST 的大小关系。把这个特例继续拆开看：节点向父亲返回的量和全局答案通常不是同一个量。先写叶子或空节点的返回值，再写父节点如何合并左右孩子，递归式才有边界基础。

\textbf{状态复用。} 把对子树的重复扫描压成返回值或队列状态；每个节点只合并孩子给出的稳定信息。本题落点是：规律是后序返回是否找到目标或祖先，不能用 BST 的大小关系。手算时只改被新输入影响的那一小部分，而不是回头重算完整候选。

\textbf{指针和字段。} 状态字段要能说清：节点指针、传入约束或返回值、当前答案、层号或路径状态；空节点的返回值必须和状态定义匹配。手算时按这个协议跟踪：拿普通二叉树根 3，p=5、q=1 手算，左右子树分别找到一个目标，根 3 就是最近公共祖先；若两个目标都在左子树，答案应由左子树返回。然后按协议复查：手算时从叶子开始写返回值，或者从根开始写传入约束。不要只写访问顺序，要写每条边上传递的信息。

\textbf{代码落点。} 递归版本先处理空节点；迭代版本的栈元素要带完整状态；全负值、重复值和极端深树要单独测。关键代码行必须能逐句翻译成“旧状态怎样变成新状态”，否则就只是模板。

\textbf{正确性。} 证明从子树归纳或层序不变量开始：孩子状态正确时父节点合并正确；若用队列，当前轮队列必须正好保存当前层节点。

\textbf{边界实验。} 先用目标不存在、p 等于 q、一个是另一个祖先、比较节点地址做反例检查。实验时覆盖目标不存在、p 等于 q、一个是另一个祖先、比较节点地址，尤其关注空节点、单节点、链式树和全负值。递归版本和显式栈版本可以互相对拍。

\textbf{复杂度变化。} 暴力在每个节点重复扫描子树或反复寻找层边界可能退化到平方级；后序汇总、显式栈或层序遍历让每个节点处理常数次，通常是线性时间。

\textbf{迁移追问。} 如果题目改成“多次 LCA 查询可预处理倍增祖先表”，先判断原状态是否仍然覆盖新答案集合，再决定是微调边界、改 value 语义，还是换算法。

\subsection{LeetCode 994 腐烂的橘子}

\textbf{题目说明。} LeetCode 994 腐烂的橘子 的题面入口要先还原成具体任务：所有初始腐烂橘子同时作为 BFS 源点。

\textbf{样例入口。} 拿 [[2, 1, 1], [1, 1, 0], [0, 1, 1]] 手算，所有腐烂橘子同时向四周扩散，第一分钟会感染相邻新鲜橘子。

\textbf{暴力基线。} 暴力按分钟模拟整张网格：每一分钟扫描所有格子，找到会被感染的新鲜橘子并统一变腐烂。

\textbf{暴力过程。} 在 [[2, 1, 1], [1, 1, 0], [0, 1, 1]] 中，初始腐烂橘子同时向四周扩散；第一分钟感染它上下左右的新鲜橘子。

\textbf{重复工作。} 题面模型：所有初始腐烂橘子同时作为 BFS 源点。暴力版的慢点要落到本题：浪费在每分钟都扫描全图。扩散只会从上一分钟刚腐烂的边界继续向外，队列正好保存这个边界。后面的优化只消掉这类重复工作，不能改变答案集合。

\textbf{优化条件。} 本题的关键观察是：按层扩散，分钟数等于最后一个新鲜橘子被首次访问的层数。这句话必须能翻译成可维护状态；如果题目条件改变后观察不再成立，就不能继续套原来的写法。

\textbf{相邻状态。} 规律是多源 BFS，初始所有腐烂橘子同时入队，层数就是分钟数。把这个特例继续拆开看：状态节点不一定等于原始位置，还可能包含钥匙、剩余资源、时间或访问集合。visited 只能合并未来能力完全相同的状态，否则会把合法路径剪掉。

\textbf{状态复用。} 多源 BFS。初始把所有腐烂橘子入队，并统计 fresh 数量；每弹出一层代表一分钟，感染相邻新鲜橘子并 fresh--。手算时只改被新输入影响的那一小部分，而不是回头重算完整候选。

\textbf{指针和字段。} 状态字段要能说清：queue 保存当前扩散边界，fresh 是剩余新鲜橘子数，minutes 是已经完成的层数，directions 是四方向。手算时按这个协议跟踪：第一层处理初始腐烂源，感染邻居入队；第二层只处理第一分钟新腐烂的橘子，层数自然对应分钟。

\textbf{代码落点。} 关键是所有初始腐烂橘子一次性入队。感染时立即把 grid[nr][nc] 改成 2，避免同一分钟被重复入队。关键代码行必须能逐句翻译成“旧状态怎样变成新状态”，否则就只是模板。

\textbf{正确性。} BFS 按层扩展，某个新鲜橘子第一次被访问的层数就是它能被腐烂的最早分钟；多源初始化表示所有源同时开始。

\textbf{边界实验。} 先用没有新鲜橘子、没有腐烂源、隔离新鲜橘子、空格阻断做反例检查。实验时覆盖没有新鲜橘子、没有腐烂源、隔离新鲜橘子、空格阻断，并打印入队时的完整状态。若同一位置但资源不同的状态被合并，很多图题会出现隐蔽错误。

\textbf{复杂度变化。} 逐分钟全图扫描最坏 O(mn*分钟数)，BFS 每格最多入队一次，O(mn) 时间、O(mn) 空间。

\textbf{迁移追问。} 如果题目改成“多源 BFS 也适用于最近零距离、地图扩散等题”，先判断原状态是否仍然覆盖新答案集合，再决定是微调边界、改 value 语义，还是换算法。

\subsection{LeetCode 210 课程表 II}

\textbf{题目说明。} LeetCode 210 课程表 II 的题面入口要先还原成具体任务：在判环基础上还要输出一个可行拓扑序。

\textbf{样例入口。} 拿 0->1、0->2 手算，课程 0 入度为 0，先输出 0，再把 1 和 2 的入度减到 0。若存在环，输出数量会少于课程数。

\textbf{暴力基线。} 慢版本先按题意画出完整状态图，用普通搜索跑通可达性。这个过程用于确认在判环基础上还要输出一个可行拓扑序的节点和边，之后再讨论访问标记、层次或代价顺序。放到本题就是：先按“在判环基础上还要输出一个可行拓扑序”完整试慢版本；样例里已经能看到“规律是课程表 II 比判定题多了输出顺序，队列弹出的顺序就是一种可行拓扑序”，所以优化前必须先能说清慢版本枚举了哪些候选。

\textbf{暴力过程。} 拿 0->1、0->2 手算，课程 0 入度为 0，先输出 0，再把 1 和 2 的入度减到 0。若存在环，输出数量会少于课程数。

\textbf{重复工作。} 题面模型：在判环基础上还要输出一个可行拓扑序。暴力版的慢点要落到本题：慢版本会围绕“在判环基础上还要输出一个可行拓扑序”反复展开同一批状态。样例规律是“规律是课程表 II 比判定题多了输出顺序，队列弹出的顺序就是一种可行拓扑序”，说明必须把节点、边、访问标记和资源维度一次建清；同一状态不应重复入队，不同资源状态也不能误合并。后面的优化只消掉这类重复工作，不能改变答案集合。

\textbf{优化条件。} 本题的关键观察是：每弹出一个入度为零课程，就把它加入顺序并删除出边。这句话必须能翻译成可维护状态；如果题目条件改变后观察不再成立，就不能继续套原来的写法。

\textbf{相邻状态。} 规律是课程表 II 比判定题多了输出顺序，队列弹出的顺序就是一种可行拓扑序。把这个特例继续拆开看：状态节点不一定等于原始位置，还可能包含钥匙、剩余资源、时间或访问集合。visited 只能合并未来能力完全相同的状态，否则会把合法路径剪掉。

\textbf{状态复用。} 把重复搜索压成访问标记、距离数组或拓扑状态；邻居生成也要从全量扫描变成结构化枚举。本题落点是：规律是课程表 II 比判定题多了输出顺序，队列弹出的顺序就是一种可行拓扑序。手算时只改被新输入影响的那一小部分，而不是回头重算完整候选。

\textbf{指针和字段。} 状态字段要能说清：状态编号、邻接生成方式、访问标记、距离或层数、队列内容；若状态含资源，visited 不能只按位置记录。手算时按这个协议跟踪：拿 0->1、0->2 手算，课程 0 入度为 0，先输出 0，再把 1 和 2 的入度减到 0。若存在环，输出数量会少于课程数。然后按协议复查：手算时画队列或栈，状态必须包含题目需要的所有资源。入队时标记还是出队时标记，要和去重语义一致。

\textbf{代码落点。} 入队时标记还是出队时标记必须一致；方向数组集中定义；多源 BFS 要一次性把所有源点放入初始队列。关键代码行必须能逐句翻译成“旧状态怎样变成新状态”，否则就只是模板。

\textbf{正确性。} 证明从可达性开始：搜索覆盖所有合法边能到达的状态，访问标记只删除重复状态而不删除不同未来能力。

\textbf{边界实验。} 先用有环返回空数组、多个可行顺序、边方向做反例检查。实验时覆盖有环返回空数组、多个可行顺序、边方向，并打印入队时的完整状态。若同一位置但资源不同的状态被合并，很多图题会出现隐蔽错误。

\textbf{复杂度变化。} 暴力重复从多个起点搜索会反复遍历图；正确建模和访问标记后，BFS 或 DFS 通常按节点数加边数计算，状态图题则按完整状态数计算。

\textbf{迁移追问。} 如果题目改成“若需要字典序最小拓扑序，把队列换成小顶堆”，先判断原状态是否仍然覆盖新答案集合，再决定是微调边界、改 value 语义，还是换算法。

\begin{principlebox}{本节复盘方式}
不要只看最终算法名。每道题复盘时先遮住优化版，口述暴力枚举对象；再指出相邻窗口、历史前缀、候选区间、搜索状态或子问题里哪一部分被重复处理；最后用一个边界样例验证状态更新顺序。能讲完这三步，才算真正掌握本章案例。
\end{principlebox}
